%==============================================================================
%  MetaLab Suite --- Problem #3
%  Population-based Multimodal Continuous Optimization: assignment (theory)
%
%  Companion technical guide: multimodal_continuous_optimization/README.md
%  Compile with:  pdflatex multimodal_continuous_optimization.tex  (twice)
%==============================================================================
\documentclass[11pt,a4paper]{article}

\usepackage[T1]{fontenc}
\usepackage[utf8]{inputenc}
\usepackage{lmodern}
\usepackage[margin=2.5cm]{geometry}
\usepackage{amsmath,amssymb,amsthm,mathtools}
\usepackage{booktabs}
\usepackage{array}
\usepackage{float}
\usepackage{xcolor}
\usepackage{microtype}
\usepackage[hidelinks]{hyperref}

\newtheorem{remark}{Remark}
\newcommand{\R}{\mathbb{R}}
\newcommand{\norm}[1]{\left\lVert #1 \right\rVert}
\DeclareMathOperator*{\argmin}{arg\,min}
\DeclareMathOperator*{\argmax}{arg\,max}

\title{%
  \textbf{Population-based Multimodal Continuous Optimization}\\[2mm]
  \large Assignment --- Theoretical Base}
\author{%
  MetaLab Suite\\[1mm]
  \small Problem \#3}
\date{\today}

%==============================================================================
\begin{document}
\maketitle

\begin{abstract}
\noindent
A swarm of sensors must triangulate the \textbf{strongest rescue signal} over a
continuous, ``n``-dimensional terrain riddled with \textbf{false positives}
(local optima). The terrain is a classic multimodal benchmark function, shifted
so that its single global optimum --- the real signal --- sits at a random
location. The scarce resource is \emph{not time}: it is the
\textbf{evaluation budget}. Every call to the terrain consumes one evaluation;
the environment counts them and cuts (or penalises) any call beyond the limit.
This document defines the problem; the technical instructions for running the
environment, the black-box solver interface and the autograder are in the
companion guide \texttt{multimodal\_continuous\_optimization/README.md}.
\end{abstract}

\tableofcontents

%==============================================================================
\section{Assignment brief}
%==============================================================================

The sensors can query the terrain at any point and read a scalar value: the
\textbf{terrain cost}. They never see the mathematical function, its name, its
gradient or the location of the signal --- only the value they measure. The goal
is to report, as accurately as possible, the \textbf{coordinates of the strongest
signal} (the global optimum), using \textbf{at most} a fixed number of
evaluations.

You must implement \textbf{one population-based metaheuristic} (Differential
Evolution, Particle Swarm Optimization, a Genetic Algorithm, CMA-ES, Artificial
Bee Colony, \dots) that operates through the black-box callback alone. The only
algorithm supplied is a plain random search.

%==============================================================================
\section{The search space}
%==============================================================================

The domain is the axis-aligned box
\begin{equation}
  \mathcal{X} = [L, U]^{n} \subset \R^{n},
  \qquad L < U,
  \label{eq:domain}
\end{equation}
where the dimension $n$ and the bounds $(L, U)$ depend on the instance (e.g.\
$n = 2$ with $[L,U] = [-5.12, 5.12]$ for the \emph{toy} terrain, $n = 10$ with
$[-500, 500]$ for the \emph{hard} one). A \emph{candidate} is a point
$x \in \mathcal{X}$; the swarm must decide which points to query.

%==============================================================================
\section{The terrain: a family of classic multimodal functions}
%==============================================================================

The terrain is one member $f : \R^{n} \to \R$ of a family of \textbf{classic
multimodal benchmark functions}. Each has a single global minimum surrounded by
many local minima --- the false positives that lure a myopic search. Three
representatives are
\begin{align}
  f_{\text{ras}}(x) &= 10n + \sum_{i=1}^{n}\Bigl(x_{i}^{2}
      - 10\cos(2\pi x_{i})\Bigr),
      && \text{(Rastrigin)}\\[1mm]
  f_{\text{ack}}(x) &= -20\exp\!\Bigl(-0.2\sqrt{\tfrac1n\textstyle\sum_i x_i^2}\Bigr)
      - \exp\!\Bigl(\tfrac1n\textstyle\sum_i\cos(2\pi x_i)\Bigr) + 20 + e,
      && \text{(Ackley)}\\[1mm]
  f_{\text{sch}}(x) &= 418.9829\,n - \sum_{i=1}^{n} x_{i}\sin\!\bigl(\sqrt{|x_{i}|}\bigr).
      && \text{(Schwefel)}
\end{align}
Table~\ref{tab:functions} lists the family the instances are drawn from, with the
canonical bounds and the \emph{canonical} optimum $x^{\mathrm{c}}$ (the argmin in
the function's own frame) and its value $f(x^{\mathrm{c}})$.

\begin{table}[H]
\centering
\caption{The multimodal family used to build terrains.}
\label{tab:functions}
\begin{tabular}{@{}llcc@{}}
\toprule
Function & Domain & Canonical argmin $x^{\mathrm{c}}$ & $f(x^{\mathrm{c}})$ \\
\midrule
Rastrigin          & $[-5.12, 5.12]$      & $0$        & $0$ \\
Ackley             & $[-32.768, 32.768]$  & $0$        & $0$ \\
Griewank           & $[-600, 600]$        & $0$        & $0$ \\
Schwefel           & $[-500, 500]$        & $420.9687$ & $0$ \\
Levy               & $[-10, 10]$          & $1$        & $0$ \\
Styblinski--Tang   & $[-5, 5]$            & $-2.9035$  & $-39.166\,n$ \\
Drop-Wave (2-D)    & $[-5.12, 5.12]$      & $0$        & $-1$ \\
Schaffer N.2 (2-D) & $[-100, 100]$        & $0$        & $0$ \\
\bottomrule
\end{tabular}
\end{table}

%==============================================================================
\section{The hidden signal: shifting the terrain}
%==============================================================================

To hide the answer, the terrain is \emph{shifted} so that its global optimum
lands at a random location $x^{\star}$ (the \textbf{signal}). Given the canonical
argmin $x^{\mathrm{c}}$, the environment evaluates
\begin{equation}
  g(x) = f\bigl(x - x^{\star} + x^{\mathrm{c}}\bigr),
  \qquad\text{so that}\qquad
  g(x^{\star}) = f(x^{\mathrm{c}}) = \min_{x \in \mathcal{X}} g(x),
  \label{eq:shift}
\end{equation}
where $x^{\star}$ is drawn uniformly from the interior of the box (kept a margin
away from the boundary). A candidate is scored by $g$, and the coordinates the
swarm reports are compared against the hidden $x^{\star}$.

\begin{remark}[Domain-limited functions]
A few functions are bounded below only on their canonical box --- most notably
\textbf{Schwefel}, which diverges to $-\infty$ outside $[-500, 500]$. Such
functions are \emph{not} shifted ($x^{\star} = x^{\mathrm{c}}$); otherwise the
``global optimum'' would drift to the boundary. The generator enforces this, so
every instance is well posed: $g(x^{\star})$ is always the true minimum.
\end{remark}

%==============================================================================
\section{The objective}
%==============================================================================

``Strongest signal'' means \emph{lowest terrain cost}: the environment returns
$g(x)$, and the swarm must \textbf{minimise} it. Equivalently, the signal
strength is $-g(x)$ and the false positives are the local minima of $g$. The
quality of a run is measured by the \textbf{distance} between the reported
coordinates and the hidden signal,
\begin{equation}
  \text{error} = \frac{\norm{\hat{x} - x^{\star}}}{\sqrt{n}\,(U - L)} \in [0, 1],
  \label{eq:error}
\end{equation}
normalised so that $0$ is a perfect hit and $1$ is the worst point in the box.

%==============================================================================
\section{The constraint: the evaluation budget}
%==============================================================================

Unlike a simulation-time problem, there is \emph{no clock} inside the model. The
single scarce resource is the number of evaluations
\begin{equation}
  \boxed{\; \sum_{k} \mathbf{1}\{\text{the swarm queries the terrain}\} \;\le\; B \;}
  \label{eq:budget}
\end{equation}
for an instance-dependent budget $B$ (e.g.\ $B = 2\,000$ for \emph{toy},
$B = 100\,000$ for \emph{hard}). The environment is the referee:

\begin{itemize}
  \item every call to the terrain \textbf{counts} one evaluation;
  \item once the budget is spent, a further call is \textbf{cut} (the environment
        raises and the run is aborted) or \textbf{penalised} (a large constant is
        returned), depending on the mode;
  \item either way the run is flagged \textbf{invalid}, so exceeding the budget
        can never beat an honest run.
\end{itemize}

\begin{remark}[Why this is the hard part]
A large budget rewards methods that \emph{use it well}, not those that iterate
most: a population that balances exploration (to leave a false positive) and
exploitation (to refine the signal) will triangulate the optimum long before a
random sampler or a trajectory method would.
\end{remark}

%==============================================================================
\section{What you are given}
%==============================================================================

The following are provided (see the companion guide for the exact commands):

\begin{itemize}
  \item the \textbf{environment engine} (\texttt{given/simulator\_mmco.py}),
  which builds the terrain, exposes the black-box evaluation and enforces the
  budget;
  \item the \textbf{instance generator}
  (\texttt{given/generate\_instances\_mmco.py}), which builds the \texttt{toy}
  and \texttt{hard} terrains and guarantees they are well posed;
  \item a \textbf{black-box student API} (\texttt{student\_mmco.py}) whose only
  window onto the terrain is the callback \texttt{evaluate(coords)};
  \item the \textbf{only algorithm supplied}, a random search
  (\texttt{random\_search\_mmco.py});
  \item a \textbf{terminal test harness} (\texttt{run\_local.py}) and the
  \textbf{autograder backend} (\texttt{autograder/}).
\end{itemize}

%==============================================================================
\section{Deliverables}
%==============================================================================

\begin{enumerate}
  \item A population-based metaheuristic that finds the signal \textbf{within the
  budget}, using only the black-box callback.
  \item Results for the \texttt{toy} and \texttt{hard} terrains: best coordinates,
  the error~\eqref{eq:error}, the evaluations used and the convergence curve.
  \item A short discussion of the exploration/exploitation balance under a fixed
  evaluation budget.
\end{enumerate}

%==============================================================================
\begin{thebibliography}{9}

\bibitem{ackley1987}
D.~H.~Ackley,
\emph{A Connectionist Machine for Genetic Hillclimbing}.
Boston, MA: Kluwer Academic Publishers, 1987.

\bibitem{kennedy1995}
J.~Kennedy and R.~Eberhart,
``Particle swarm optimization,''
in \emph{Proc.\ IEEE International Conference on Neural Networks}, 1995,
pp.~1942--1948.

\bibitem{storn1997}
R.~Storn and K.~Price,
``Differential evolution --- a simple and efficient heuristic for global
optimization over continuous spaces,''
\emph{Journal of Global Optimization}, vol.~11, no.~4, pp.~341--359, 1997.

\bibitem{karaboga2007}
D.~Karaboga and B.~Basturk,
``A powerful and efficient algorithm for numerical function optimization:
artificial bee colony (ABC) algorithm,''
\emph{Journal of Global Optimization}, vol.~39, no.~3, pp.~459--471, 2007.

\bibitem{hansen2001}
N.~Hansen and A.~Ostermeier,
``Completely derandomized self-adaptation in evolution strategies,''
\emph{Evolutionary Computation}, vol.~9, no.~2, pp.~159--195, 2001.

\bibitem{talbi2009}
E.-G.~Talbi,
\emph{Metaheuristics: From Design to Implementation}.
Hoboken, NJ: John Wiley \& Sons, 2009.

\end{thebibliography}

\end{document}
